
\documentclass{article}
\usepackage{colortbl}
\usepackage{makecell}
\usepackage{multirow}
\usepackage{supertabular}

\begin{document}

\newcounter{utterance}

\centering \large Interaction Transcript for game `cladder', experiment `full\_v1.5\_default', episode 9125 with llama{-}3.1{-}8b{-}instruct.
\vspace{24pt}

{ \footnotesize  \setcounter{utterance}{1}
\setlength{\tabcolsep}{0pt}
\begin{supertabular}{c@{$\;$}|p{.15\linewidth}@{}p{.15\linewidth}p{.15\linewidth}p{.15\linewidth}p{.15\linewidth}p{.15\linewidth}}
    \# & \multicolumn{2}{c}{Player} && \multicolumn{2}{c}{Game Master} \\
    \hline

    \theutterance \stepcounter{utterance}  
    & & & \multicolumn{4}{p{0.6\linewidth}}{
        \cellcolor[rgb]{0.9,0.9,0.9}{
            \makecell[{{p{\linewidth}}}]{
                \texttt{\tiny{[P1$\langle$GM]}}
                \texttt{You are an expert in causal inference. The following question is not a typical commonsense query, but rather a meticulously designed question created by a professor specializing in causal inference, intended to assess the students' mastery of the course content.} \\
\texttt{Imagine a self{-}contained, hypothetical world with only the following conditions, and without any unmentioned factors or causal relationships: Kraz has a direct effect on pexu and rukz. Pexu has a direct effect on hwax. Hwax has a direct effect on rukz. Kraz is unobserved. For those who are not pexu and are not hwax, the probability of rukz is 19\%. For those who are not pexu and are hwax, the probability of rukz is 29\%. For those who are pexu and are not hwax, the probability of rukz is 20\%. For those who are pexu and are hwax, the probability of rukz is 31\%. For those who are not pexu, the probability of hwax is 16\%. For those who are pexu, the probability of hwax is 90\%. For those who are pexu, would it be less likely to see rukz if the individual was not pexu?} \\
\texttt{Start your answer with 'Yes' or 'No', followed by additional reasoning or evidence to support your explanation.} \\
            }
        }
    }
     \\ \\

    \theutterance \stepcounter{utterance}  
    & \multicolumn{4}{p{0.6\linewidth}}{
        \cellcolor[rgb]{0.561,0.737,0.561}{
            \makecell[{{p{\linewidth}}}]{
                \texttt{\tiny{[P1$\rangle$GM]}}
                \texttt{No.} \\
\\ 
\texttt{To address this question, we need to apply the principles of causal inference, specifically the concept of counterfactuals and potential outcomes. } \\
\\ 
\texttt{Let's denote the potential outcome of rukz if an individual is not pexu as Rukz(0), and the potential outcome of rukz if an individual is pexu as Rukz(1). We are interested in comparing Rukz(0) and Rukz(1) for individuals who are pexu.} \\
\\ 
\texttt{We can use the law of total probability to calculate the probability of rukz for individuals who are pexu:} \\
\\ 
\texttt{P(Rukz|Pexu) = P(Rukz|Pexu, Hwax) * P(Hwax|Pexu) + P(Rukz|Pexu, \textasciitilde{}Hwax) * P(\textasciitilde{}Hwax|Pexu)} \\
\\ 
\texttt{Given the provided probabilities, we can substitute the values:} \\
\\ 
\texttt{P(Rukz|Pexu) = 0.31 * 0.9 + 0.20 * 0.1} \\
\texttt{P(Rukz|Pexu) = 0.279 + 0.02} \\
\texttt{P(Rukz|Pexu) = 0.299} \\
\\ 
\texttt{Now, we need to compare this probability with the probability of rukz for individuals who are not pexu but are in the same stratum (i.e., not pexu and hwax). We are given that P(Rukz|\textasciitilde{}Pexu, Hwax) = 0.29.} \\
\\ 
\texttt{To compare the two probabilities, we need to consider the probability of being in the same stratum (not pexu and hwax) for individuals who are pexu. However, this probability is not directly provided. } \\
\\ 
\texttt{Instead, we can use the law of total probability to calculate the probability of being in the same stratum (not pexu and hwax) for individuals who are pexu:} \\
\\ 
\texttt{P(\textasciitilde{}Pexu, Hwax|Pexu) = P(\textasciitilde{}Pexu|Pexu) * P(Hwax|\textasciitilde{}Pexu, Pexu)} \\
\\ 
\texttt{Given the provided probabilities, we can substitute the values:} \\
\\ 
\texttt{P(\textasciitilde{}Pexu|Pexu) = 1 {-} P(Pexu|Pexu) = 1 {-} 1 = 0 (this is incorrect, we should use the given probability of hwax for pexu individuals, which is 0.9, to calculate the probability of not pexu for pexu individuals. However, we can use the given probability of hwax for pexu individuals to calculate the probability of not hwax for pexu individuals, which is 0.1)} \\
\\ 
\texttt{P(Hwax|\textasciitilde{}Pexu, Pexu) = P(Hwax|\textasciitilde{}Pexu) / P(\textasciitilde{}Pex} \\
            }
        }
    }
    & & \\ \\

    \theutterance \stepcounter{utterance}  
    & & & \multicolumn{2}{p{0.3\linewidth}}{
        \cellcolor[rgb]{0.9,0.9,0.9}{
            \makecell[{{p{\linewidth}}}]{
                \texttt{\tiny{[GM$|$GM]}}
                \texttt{No} \\
            }
        }
    }
    & & \\ \\

    \theutterance \stepcounter{utterance}  
    & & & \multicolumn{2}{p{0.3\linewidth}}{
        \cellcolor[rgb]{0.9,0.9,0.9}{
            \makecell[{{p{\linewidth}}}]{
                \texttt{\tiny{[GM$|$GM]}}
                \texttt{yes} \\
            }
        }
    }
    & & \\ \\

    \theutterance \stepcounter{utterance}  
    & & & \multicolumn{2}{p{0.3\linewidth}}{
        \cellcolor[rgb]{0.9,0.9,0.9}{
            \makecell[{{p{\linewidth}}}]{
                \texttt{\tiny{[GM$|$GM]}}
                \texttt{game\_result = LOSE} \\
            }
        }
    }
    & & \\ \\

\end{supertabular}
}

\end{document}
